\documentclass[11pt]{article}

\usepackage[utf8]{inputenc}
\usepackage[T1]{fontenc}
\usepackage[serbian]{babel}
\usepackage[margin=1.9cm]{geometry}
\usepackage{amsmath,amssymb}
\usepackage{algorithm}
\usepackage{algpseudocode}
\usepackage{enumitem}
\usepackage{amsthm}
\usepackage{hyperref}
\usepackage{xcolor}
\usepackage{titlesec}
\usepackage{booktabs}

\titlespacing*{\section}{0pt}{5pt}{2pt}
\titlespacing*{\subsection}{0pt}{3pt}{1pt}
\setlength{\parskip}{1pt}
\setlength{\abovedisplayskip}{2pt}
\setlength{\belowdisplayskip}{2pt}
\setlength{\textfloatsep}{3pt}
\setlength{\floatsep}{2pt}
\setlength{\intextsep}{2pt}
\renewcommand{\baselinestretch}{0.96}

\theoremstyle{definition}
\newtheorem{definition}{Definicija}[section]
\newtheorem{theorem}{Teorema}[section]
\floatname{algorithm}{Algoritam}

\hypersetup{colorlinks=true, linkcolor=blue, urlcolor=blue, citecolor=blue}

\begin{document}
\sloppy

\title{\vspace{-1.5cm}Problem \emph{Polyline Simplification}
\\[0.2cm]
\small{Projekat iz predmeta Geometrijski algoritmi\\ Matematički fakultet, Univerzitet u Beogradu}}
\author{Boško Andrić, 1012/2025\\ \texttt{mi251012@alas.matf.bg.ac.rs}}
\date{\today}
\maketitle

\vspace{-0.5cm}
\tableofcontents
\vspace{0.3cm}

\section{Uvod i opis problema}

Smanjivanje broja tačaka na krivim linijama je čest zadatak u kartografiji i obradi geografskih podataka. Problem \emph{Polyline Simplification} sa Kattis platforme \cite{kattis} traži da se iz izlomljene linije date nizom tačaka $P = (p_0, p_1, \ldots, p_n)$ izbace manje važne tačke. Cilj je da od početnih $n+1$ temena ostane tačno $m$ segmenata, a da linija zadrži svoj prepoznatljiv oblik.

U ovom projektu primenjen je Visvalingam--Whyatt algoritam \cite{visvalingam1993}. Važnost svake unutrašnje tačke računa se na osnovu površine trougla koji ona pravi sa svojim susedima. U svakom koraku izbacuje se tačka sa najmanjom površinom trougla. Ako više tačaka ima istu površinu, prva se izbacuje ona sa manjim indeksnim brojem.

Ako bismo u svakom koraku tražili najmanju površinu prolaskom kroz sve tačke, složenost bi bila $O(n^2)$. Korišćenjem pametnijih struktura podataka, taj posao smanjujemo na vremensku složenost $O(n \log n)$ uz memorijsku složenost $O(n)$.

\section{Geometrijski princip}

\subsection{Dvostruka površina trougla}

Za unutrašnju tačku $p_i = (x_i, y_i)$ posmatramo trougao koji ona čini sa susedima $p_{i-1}$ i $p_{i+1}$.

\begin{definition}[Površina trougla]
\label{def:povrsina}
Važnost tačke $p_i$ jednaka je površini trougla $\triangle(p_{i-1}, p_i, p_{i+1})$. Da bismo izbegli rad sa razlomcima i deljenjem, računamo dvostruku površinu pomoću vektorskog proizvoda:
\[
\mathcal{A}(p_{i-1}, p_i, p_{i+1}) = \big| (x_i - x_{i-1})(y_{i+1} - y_{i-1}) - (y_i - y_{i-1})(x_{i+1} - x_{i-1}) \big|
\]
\end{definition}

Sva poređenja radimo direktno preko ovih celobrojnih vrednosti $\mathcal{A}$, što ubrzava algoritam i izbegava greške pri zaokruživanju.

\subsection{Lokalni efekat izbacivanja}

Efikasnost algoritma oslanja se na činjenicu da brisanje jedne tačke utiče samo na njene prve susede:

\begin{theorem}[Lokalno osvežavanje]
\label{thr:lokalnost}
Uklanjanje tačke $p_i$ menja trouglove samo za njenog prethodnika $p_{\text{prev}(i)}$ i njenog sledbenika $p_{\text{next}(i)}$. Površine ostalih trouglova ostaju potpuno iste.
\end{theorem}

\begin{proof}
Kada obrišemo tačku $p_i$, spajaju se njeni susedi $p_L$ i $p_R$. Zbog toga se menja trougao za $p_L$ i trougao za $p_R$. Pošto se ostale tačke i njihove veze nisu menjale, njihove površine ostaju nepromenjene.
\end{proof}

Zahvaljujući ovome, posle svakog brisanja ažuriramo površine za samo dve susedne tačke u $O(\log n)$ vremenu, umesto da računamo sve ponovo.

\section{Struktura programa i algoritam}

Program je napisan u C++17 jeziku koristeći Qt 6 za prikaz. Logika algoritma nalazi se u \texttt{src/efficientsimplifier.cpp}, dok se prikaz i prozor nalaze u \texttt{src/canvas.cpp} i \texttt{src/mainwindow.cpp}.

\subsection{Strukture podataka za $O(n \log n)$ složenost}

Za postizanje $O(n \log n)$ složenosti upotrebljene su dve strukture:
\begin{enumerate}[leftmargin=*]
    \item \textbf{Dvostruko povezana lista (preko nizova):} Nizovi \texttt{prev\_neighbor} i \texttt{next\_neighbor} služe da u $O(1)$ vremenu nađemo i prevežemo susede kada se tačka obriše.
    \item \textbf{Prioritetni skup (\texttt{std::set}):} Čuva indekse tačaka i njihove površine. Skup je sortiran primarno po površini (od najmanje ka najvećoj), a sekundarno po indeksu tačke.
\end{enumerate}

\subsection{Koraci algoritma}

Algoritam~\ref{alg:polyline} prikazuje jednostavan tok simplifikacije.

\begin{algorithm}[h]
\caption{Simplifikacija polilinije u $O(n \log n)$ vremenu}
\label{alg:polyline}
\begin{algorithmic}[1]
\Function{VisvalingamSimplify}{$P = (p_0, \ldots, p_n)$, $m$}
    \State Inicijalizuj nizove \texttt{prev\_neighbor} i \texttt{next\_neighbor}
    \State $Q \gets \emptyset$ \Comment{Prioritetni skup \texttt{std::set}}
    \For{$i \gets 1$ \textbf{to} $n-1$}
        \State $\text{area}[i] \gets \text{ComputeDoubleArea}(p_{i-1}, p_i, p_{i+1})$
        \State $Q.\text{insert}(\{i, \text{area}[i]\})$
    \EndFor
    \State $count \gets n+1$
    \While{$count > m+1$ \textbf{and} $Q \neq \emptyset$}
        \State $min\_elem \gets Q.\text{extract\_min}()$ \Comment{Uzimamo najmanju površinu u $O(\log n)$}
        \State $idx \gets min\_elem.index$, $L \gets \text{prev}[idx]$, $R \gets \text{next}[idx]$
        \State $\text{next}[L] \gets R$, $\text{prev}[R] \gets L$ \Comment{Prevezivanje suseda u $O(1)$}
        \State Sačuvaj korak u istoriju za animaciju
        \If{$L > 0$} \Comment{Ažuriraj levog suseda}
            \State $Q.\text{erase}(\{L, \text{area}[L]\})$
            \State $\text{area}[L] \gets \text{ComputeDoubleArea}(p_{\text{prev}[L]}, p_L, p_R)$
            \State $Q.\text{insert}(\{L, \text{area}[L]\})$
        \EndIf
        \If{$R < n$} \Comment{Ažuriraj desnog suseda}
            \State $Q.\text{erase}(\{R, \text{area}[R]\})$
            \State $\text{area}[R] \gets \text{ComputeDoubleArea}(p_L, p_R, p_{\text{next}[R]})$
            \State $Q.\text{insert}(\{R, \text{area}[R]\})$
        \EndIf
        \State $count \gets count - 1$
    \EndWhile
    \State \Return istorija koraka
\EndFunction
\end{algorithmic}
\end{algorithm}

\section{Složenost}

\begin{center}
\begin{tabular}{@{}lll@{}}
\toprule
\textbf{Korak} & \textbf{Složenost} & \textbf{Opis} \\
\midrule
Inicijalizacija i prve površine & $O(n)$ & Prolazak kroz sve tačke jednim nizom \\
Pravljenje \texttt{std::set} stablom & $O(n \log n)$ & Ubacivanje $n-1$ elemenata u skup \\
Uzimanje najmanje površine & $O(\log n)$ & Brisanje prvog elementa iz skupa \\
Ažuriranje 2 susedne tačke & $O(\log n)$ & Izbacivanje i ponovno ubacivanje u skup \\
Ukupno za sve izbačene tačke & $O(n \log n)$ & Mnogo brže od naivnog $O(n^2)$ pristupa \\
\midrule
\textbf{Utoršak memorije} & $O(n)$ & Čuvanje nizova suseda i stanja u skupu \\
\bottomrule
\end{tabular}
\end{center}

Naivni algoritam za $n = 100.000$ tačaka pravi oko $10^{10}$ računica, dok ova optimizovana verzija završava posao za manje od sekunde.

\section{Prikaz u aplikaciji}

Strogi poredak u \texttt{std::set} strukturi garantuje da će se uvek izbaciti tačna tačka propisana algoritmom.

Aplikacija podržava dva načina prikaza:
\begin{itemize}[leftmargin=*]
    \item \textbf{Algoritamski mod:} Prikazuje originalnu liniju sivom isprekidanom linijom, trenutnu liniju plavom bojom, a trougao koji se sledeći izbacuje označava crvenom bojom.
    \item \textbf{Planinski mod (LOD):} Prikazuje liniju kao reljef planine. Kako se tačke izbacuju, slika se postepeno smanjuje i skalira, dajući efekat odaljavanja kamere.
\end{itemize}

\section{Zaključak}

U projektu je uspešno napravljen Visvalingam--Whyatt algoritam za simplifikaciju polilinija u vremenu $O(n \log n)$. Spajanjem dvostruko povezane liste i balansiranog stabla postignuta je velika brzina rada. Grafički interfejs u Qt 6 okruženju omogućava jednostavan prikaz koraka algoritma i vizuelni efekat smanjenja detalja.

\begin{thebibliography}{9}
\bibitem{kattis} Kattis, \emph{Polyline Simplification}, \url{https://open.kattis.com/problems/polyline-simplification}.
\bibitem{visvalingam1993} M.~Visvalingam, J.~D.~Whyatt, \emph{Line simplification by repeated elimination of points}, The Cartographic Journal, 30(1), 46-51, 1993.
\bibitem{janicic} P.~Janičić, \emph{Računarska geometrija}, Matematički fakultet u Beogradu, 2025.
\bibitem{cormen} T.~H.~Cormen et al., \emph{Introduction to Algorithms}, MIT Press, 2009.
\bibitem{qt6} Qt Group, \emph{Qt 6 Documentation}, \url{https://doc.qt.io/qt-6/}.
\end{thebibliography}

\end{document}